\section{An exact period on the actual linked inputs}

Retain $u,B,C,e,d$ from \eqref{eq:delta}, and write
\[
 B=2^K,\qquad r=K-1,\qquad q=2^r+1=B/2+1,\qquad
 m_0=u-2,\qquad \Delta=m_0(m_0+1)/2.
\]
The integer $q$ need not be prime. All primitive $q$-th roots
in an algebraic closure of $\Ftwo$ will be used.

\begin{lemma}[The actual residue formula]\label{lem:residue}
The form $\omega_e=z^{e-1}A_d(z)^{-1}\,dz$ has a simple nonzero
residue at each primitive $q$-th root $\zeta$, given by
\[
 b_\zeta=\frac{\zeta^e}{A_{d-q}(\zeta)}
                 =\zeta^{e-\Delta}A_{m_0}(\zeta).
\]
\end{lemma}
\begin{proof}
The hierarchy bounds give $d/2<q\le d<2q$ and $m_0<q$.
Thus only $1+z^q$ among the factors of $A_d$ vanishes at
$\zeta$, and its derivative there is $\zeta^{-1}$.
The product $\prod_{i=1}^{q-1}(1+\zeta^i)$ is one: it is
the derivative at one of $X^q-1$, whose value is $q=1$ in
characteristic two. It follows that the residue is the first
displayed expression.

Now $d-q=q-1-m_0$. The last $m_0$ terms in that complete
product equal
$\prod_{j=1}^{m_0}(1+\zeta^{-j})
  =\zeta^{-\Delta}A_{m_0}(\zeta)$.
Taking the complementary product proves the second expression.
None of $1+\zeta^j$ for $1\le j\le m_0$ vanishes.
\end{proof}

\begin{lemma}[The bounds needed for the root count]\label{lem:bounds}
For the actual linked inputs, $r\ge96$, $\Delta<r$, and
\[
 0<\Delta-2e<q,\qquad
 \varphi(q)>\frac{2^r}{r+1}>2^{2r/3}\Delta.
\]
\end{lemma}
\begin{proof}
For $s=u+100$, $F_5(s)\ge F_1(s)\ge2^{2^s}$, so
$K\ge2^{u+100}$ and $r\ge2^{u+100}-1$.
These inequalities imply $r\ge96$ and
$\Delta=(u-2)(u-1)/2<r$.
As $e=u-4=m_0-2$,
$\Delta-2e=(m_0^2-3m_0+8)/2$ is positive and smaller
than $q$ by the same growth bound.

If an odd integer $q$ has $k$ distinct prime factors
$p_1<\cdots<p_k$, then $p_i\ge i+2$, hence
$q/\varphi(q)=\prod p_i/(p_i-1)\le k+1$.
Here $3^k\le q=2^r+1$ implies $k\le r$, so
$\varphi(q)\ge q/(r+1)>2^r/(r+1)$.
Finally $2^{r/3}>r^3$ at $r=96$ and thereafter:
the ratio of the exponential to $r^3$ increases for $r\ge96$.
Since $\Delta<r$ and $r(r+1)<r^3$, this gives
$2^r/(r+1)>2^{2r/3}\Delta$.
\end{proof}

\begin{theorem}[Exact actual residue-tuple period]\label{thm:actual-period}
The tuple $(b_\zeta)_{\zeta\ {\rm primitive}\ q}$ has exact
Frobenius period $2r$. Every eventual period of the complete
rational differential Cartier orbit of $\omega_e$ is divisible
by $2r=2(\log_2 B-1)$.
\end{theorem}
\begin{proof}
The order of two modulo $q=2^r+1$ is $2r$.
It divides $2r$ because $2^r\equiv-1\pmod q$.
Any proper divisor of $2r$ is at most $r$, and for
$s\le r$ the positive integer $2^s-1$ is smaller than $q$.
Thus no smaller order is possible. The primitive roots and their
residues belong to $\mathbb F_{2^{2r}}$, so the tuple period
divides $2r$.

Reciprocity gives
$A_{m_0}(\zeta^{-1})=\zeta^{-\Delta}A_{m_0}(\zeta)$.
If Frobenius power $r$ fixed the tuple,
Lemma~\ref{lem:residue} would give
$\zeta^{\Delta-2e}=1$ at a primitive root, contrary to
$0<\Delta-2e<q$.

Any other proper divisor $s$ of $2r$ is at most $2r/3$:
the integer quotient $(2r)/s$ is then at least three.
Fixedness at all primitive roots would make every such root a zero of
\[
 P_s(X)=A_{m_0}(X)^{2^s}
       +X^{(\Delta-e)(2^s-1)}A_{m_0}(X).
\]
This polynomial has constant coefficient one and degree exactly
$2^s\Delta$, since $e>0$ makes the second summand's degree
strictly smaller. It cannot have all $\varphi(q)$ primitive roots
as zeros, by Lemma~\ref{lem:bounds}. Thus no proper divisor is a
tuple period.

The chosen finite poles were already simple initially.
By Lemma~\ref{lem:local}, the tuple in the $j$-th Cartier state is
$(b_\zeta^{\,2^{-j}})_\zeta$ with poles fixed. If states at
$j$ and $j+p$ agree, these tuples agree. Inverse Frobenius is
bijective and its tuple orbit has period $2r$, so $2r\mid p$.
This applies also after the higher-pole transient.
\end{proof}

The proof is a tuple argument. It does not assert that every individual
residue has degree $2r$, and it does not calculate the selected
constant coefficient. In particular, the growing tuple period does
not prove that the sequence $v\mapsto\delta_4(v)$ is nonperiodic.

